\section{Temperance: Cardinal Virtue and Its Parts}
\label{sec:temperance}

Temperance preserves reason's order within attractions that are good in themselves but can absorb attention and desire. In the strict sense it moderates the strongest pleasures of touch; by extension its annexed virtues moderate anger, self-estimation, desire for knowledge, bodily presentation, and play. Aristotle treats temperance in \emph{Nicomachean Ethics} III.10--12 and several lesser means in IV; Aquinas's full architecture occupies ST II--II, qq.~141--170.

\cardinalgroup{The principal virtue and the pleasures of touch}

\virtuedossier
  {Temperance}
  {Principal cardinal virtue; Aristotle, \emph{Nicomachean Ethics} II.7 and III.10--12; Augustine, \emph{On the Morals of the Catholic Church} I.15.25; ST II--II, qq.~141--145; CCC 1809.}
  {Pleasures and desires of touch, especially those attached to food, drink, and sex, measured by bodily good, personal vocation, justice, and reason.}
  {To desire and enjoy sensible goods in the manner, measure, circumstances, and end that sustain the whole human good.}
  {\textbf{Insensibility} rejects fitting bodily pleasure contrary to reason (\oppositioncode{M}).}
  {\textbf{Self-indulgence} lets pleasure overrule reason (\oppositioncode{M}).}
  {Temperance is not absence of appetite, fear of the body, aesthetic minimalism, or a naturally low intensity of desire. Integral ``shamefacedness'' and love of moral beauty dispose or ornament it but are not coordinate virtues.}

Temperance does not ask whether pleasure is intense in the abstract; it asks whether desire answers a genuinely fitting good. A feast can be temperate and a small meal gluttonous when the former expresses communion and proportion while the latter is chosen through domination by appetite. Voluntary abstention can serve health, solidarity, worship, or freedom; it becomes disordered when it rejects a creature as evil, violates duty, or seriously harms the person without proportionate reason.

\cardinalgroup{Integral parts: shamefacedness and moral beauty}

\begin{studybox}{Conditions of complete temperance}
Aquinas identifies \textbf{shamefacedness} and \textbf{honesty} or love of moral beauty as integral parts of temperance (ST II--II, qq.~143--145). Shame is a passion that recoils from disgrace and is especially useful before virtue is settled; it is treated later as a qualified state, not counted as another strict virtue. Moral beauty names the fitting radiance of ordered action and belongs broadly to virtue rather than adding a separate habit.
\end{studybox}

\cardinalgroup{Subjective species: food, drink, and sexual appetite}

\virtuedossier
  {Abstinence}
  {Subjective species of temperance; ST II--II, qq.~146--148.}
  {Food and the pleasures of eating as ordered to health, shared life, discipline, and higher goods.}
  {To take or forego food according to right quantity, kind, time, manner, and end.}
  {Unreasonable refusal of nourishment is a real disorder but has no virtue-specific source name (\oppositioncode{U}); it falls under insensibility or another failure according to its object.}
  {\textbf{Gluttony} is the named local contrary that subordinates reason to the pleasure or pursuit of food (\oppositioncode{K}); it is not one member of a complete named pair.}
  {Fasting is an act that abstinence and religion can command, not another habit. Body size, appetite, illness, poverty, culinary pleasure, and eating frequency do not by themselves reveal virtue or vice.}

\virtuedossier
  {Sobriety}
  {Subjective species of temperance; ST II--II, qq.~149--150.}
  {Intoxicating drink and its power to alter rational use. By extension the term may be applied analogically, but the controlling dossier is specific.}
  {To use or abstain from alcohol so that reason, responsibility, health, and charity retain their order.}
  {Culpable abstention that gravely harms bodily good would be vicious, but Aquinas expressly leaves it unnamed (\oppositioncode{U}); abstinence from alcohol is not by itself a failure.}
  {\textbf{Drunkenness} is the named local contrary that voluntarily disorders or surrenders rational use through intoxication (\oppositioncode{K}), with culpability depending on knowledge and freedom.}
  {Sobriety is not a diagnosis of addiction or a judgment based on one visible beverage. Medical vulnerability, scandal, law, office, and duty can make abstention the prudent act.}

\virtuedossier
  {Chastity}
  {Subjective species of temperance; Augustine, \emph{City of God} I.18, on chastity as a virtue of the soul governed by the will; ST II--II, qq.~151 and 153--154; CCC 2337.}
  {Sexual appetite and acts as integrated with the truth of the person, vocation, fidelity, generation, and charity.}
  {To govern sexual desire and conduct so that they express a just and integral personal good rather than using oneself or another.}
  {No virtue-specific defect-side vice is named (\oppositioncode{U}); insensibility, fear, trauma, or lack of desire must not be relabeled as moral vice.}
  {\textbf{Lust} is the named local contrary that seeks sexual pleasure against reason and the goods specifying the act (\oppositioncode{K}); no coordinate defect flank is supplied.}
  {Chastity is not simply non-activity, bodily intactness, repression, or one vocation. Married and unmarried chastity have different fitting acts under one virtue. ``Purity'' in ST II--II, q.~151, a.~4 is not a second coordinate virtue.}

\virtuedossier
  {Virginity}
  {Special virtue within chastity, formally grounded in a deliberate and stable purpose; Augustine, \emph{On Holy Virginity} 8, 18; ST I--II, q.~64, a.~1 ad 3; II--II, q.~152, aa.~1, 3.}
  {Perpetual renunciation of sexual intercourse for a fitting spiritual good, undertaken by a deliberate and stable purpose.}
  {To preserve that purpose and embody undivided dedication according to one's state and duties.}
  {An unnamed circumstantial deficiency fails to keep an undertaken purpose when or as due (\oppositioncode{U}); marriage and chaste marital intercourse are goods, not deficiencies of virtue.}
  {Undertaking or observing virginity for superstition, vainglory, or another undue end is a circumstantial excess without a dedicated vice proper to virginity (\oppositioncode{U}).}
  {Bodily integrity can exist without the virtue and can be lost without voluntary fault; the virtue formally lies in the purpose. Virginity is a counsel and vocation, not the maximum possible ``quantity'' of chastity.}

\cardinalgroup{Potential parts: continence, anger, and lesser matters}

Potential parts extend temperance's restraint to secondary matters. Continence resists disordered desire but lacks the completed harmony of strict moral virtue and is therefore treated among the qualified states. Clemency and meekness govern punishment and anger; modesty gathers lesser movements and contains more specifically named virtues where a distinct object warrants them (ST II--II, q.~143, a.~1; qq.~155--169).

\virtuedossier
  {Clemency}
  {Annexed virtue of temperance; ST II--II, qq.~157 and 159.}
  {A superior's impulse and authority to impose punishment upon a person subject to judgment.}
  {To mitigate penalty where reason permits while preserving correction, protection, and justice.}
  {\textbf{Laxity} is the compact attested description of undue remission: mitigating punishment against justice (\oppositioncode{U}; ST II--II, q.~159, a.~2, obj.~3). Aquinas does not establish it as a dedicated vice species.}
  {\textbf{Cruelty} inclines one to exceed rational order in punishment (\oppositioncode{M}).}
  {Clemency presupposes authority to punish and differs from meekness, which moderates one's anger, and mercy, which responds to misery. It cannot waive another person's claim without standing.}

\virtuedossier
  {Meekness}
  {Annexed virtue of temperance and Aristotelian mean often rendered gentleness; Aristotle, \emph{Nicomachean Ethics} II.7 and IV.5; ST II--II, q.~157.}
  {Anger: the passion that seeks resistance or redress under an apprehended injury.}
  {To become angry, remain calm, or act against injury according to the right cause, person, intensity, time, and manner.}
  {\textbf{Spiritlessness} fails to resist wrong when reason requires anger and action (\oppositioncode{M}).}
  {\textbf{Irascibility} lets anger exceed reason in cause, object, intensity, duration, or expression (\oppositioncode{M}).}
  {Meekness is not quiet temperament, conflict avoidance, inability to feel anger, or tolerance of abuse. Just anger remains an instrument rather than the ruler of action.}

\cardinalgroup{Modesty and its specified fields}

In Aquinas, \emph{modestia} is the common moderation of matters less difficult than temperance's principal pleasures. Some of its fields receive special virtues because their objects supply distinct rational measures: humility for the appetite of excellence, studiousness for knowledge, and eutrapelia for play. Bodily movement and dress remain applications of modesty rather than additional coordinate habits (ST II--II, qq.~160--169).

\virtuedossier
  {Modesty}
  {General annexed virtue of temperance in lesser movements; ST II--II, q.~160.}
  {The outward and inward movements not already governed by a more specifically named temperance-family virtue.}
  {To give gesture, bearing, presentation, and lesser impulses a form fitting person, place, role, culture, and end.}
  {No single dedicated defect vice covers modesty's diverse matter (\oppositioncode{U}); neglect of fitting expression must be specified by context.}
  {No single dedicated excess vice covers it (\oppositioncode{U}); theatricality, display, and rigidity are different failures rather than one canonical extreme.}
  {Modesty is not reduced to sexual conduct or clothing, and culturally unfamiliar expression is not presumptively vice. Humility, studiousness, and eutrapelia receive separate dossiers because their formal objects differ.}

\virtuedossier
  {Humility}
  {Specified virtue under modesty; Augustine, \emph{City of God} XIV.13 and \emph{On Holy Virginity} 31--32 (the latter mediated through Aquinas's citations); ST II--II, q.~161, aa.~1--2.}
  {One's appetite for excellence and standing, measured by truth about creaturely dependence, gifts, limits, vocation, and God.}
  {To restrain claims beyond one's measure and to submit oneself rightly without denying gifts entrusted for service.}
  {No distinct defect-side vice is supplied (\oppositioncode{U}). False self-lowering can be pride's counterfeit expression rather than a separate flank.}
  {\textbf{Pride} is the named contrary that seeks excellence against right order, especially through refusal of due subjection to God (\oppositioncode{K}).}
  {Humility and magnanimity are complementary: truth restrains false greatness and also forbids pusillanimous denial of real capacity. Humiliation suffered, low status, timidity, and self-dislike are not the virtue.}

\virtuedossier
  {Studiousness}
  {Specified virtue under modesty; ST II--II, qq.~166--167.}
  {The desire and effort to know, ordered by the truth sought, one's duties, method, authority, time, and the use of knowledge.}
  {To apply the mind earnestly and proportionately to knowledge worth seeking.}
  {\textbf{Neglect} of study is the defect-side failure when bodily aversion to effort keeps the mind from knowledge that should be sought (\oppositioncode{M}; ST II--II, q.~166, a.~2).}
  {\textbf{Curiosity} is the excess-side disorder that pursues knowledge in a wrong object, motive, method, priority, or use (\oppositioncode{M}).}
  {Studiousness differs from the intellectual virtues that perfect knowledge already possessed. Information volume, academic degree, novelty, and intrusive access are not measures of its excellence.}

\virtuedossier
  {Eutrapelia (ready wit in play)}
  {Specified virtue under modesty and Aristotelian mean; Aristotle, \emph{Nicomachean Ethics} II.7 and IV.8; ST II--II, q.~168, aa.~2--4.}
  {Rest, games, humor, and playful conversation as necessary recreation for a finite rational life.}
  {To give and receive fitting amusement in a way proportioned to persons, time, subject, and moral good.}
  {\textbf{Boorishness} rejects or spoils fitting play (\oppositioncode{M}).}
  {\textbf{Buffoonery} sacrifices truth, dignity, duty, or another's good to amusement (\oppositioncode{M}).}
  {Eutrapelia is not constant wit, entertainment skill, sarcasm, or exemption from reverence. Serious temperament can possess it through a fitting willingness to rest.}

\begin{studybox}{Bearing and dress}
ST II--II, qq.~168--169 applies modesty to bodily movement, voice, display, and clothing. The measure includes truthful social meaning, decency, role, custom, resources, beauty, and the avoidance of vanity or culpable negligence. No garment, luxury level, style, or gesture is virtuous in abstraction. Because Aquinas returns these applications to modesty rather than establishing new habits, the census nests them here and does not create ``decorum'' and ``dress moderation'' rows.
\end{studybox}

\cardinalgroup{Temperance as liberation}

Pleasure, anger, recognition, knowledge, and play are human goods or powers for good; temperance does not compete with their created goodness. It frees them from becoming total claims. This is why temperance can appear both restrained and abundant: it says no to domination so that appetite can say a fuller yes to health, communion, truth, worship, beauty, and love.
